% !TEX root = ../Monografia.tex
% ---------------------------------------------------------------------------------------------------------------------------------------------------------
% CAPÍTULO X - METODOLOGIA E DESENVOLVIMENTO
% ---------------------------------------------------------------------------------------------------------------------------------------------------------

\chapter{Metodologia e Desenvolvimento}


% ---------------------------------------------------------------------------------------------------------------------------------------------------------
\section{Abordagem Metodológica}
% ---------------------------------------------------------------------------------------------------------------------------------------------------------

\subsection{Natureza da Pesquisa}

A presente pesquisa caracteriza-se como um estudo de natureza \textbf{aplicada}, pois visa desenvolver e validar funcionalmente um pacote de monitoramento de desempenho para aplicações Flutter.
Segundo Gil~\cite{gil2017}, pesquisas aplicadas têm como propósito “gerar conhecimentos voltados à aplicação prática, dirigidos à solução de problemas específicos”.

Nesse contexto, o \texttt{collector\_flutter} oferece uma solução técnica para coletar, analisar e interpretar métricas de execução em tempo real, apoiando o aprimoramento da eficiência e da sustentabilidade digital de aplicativos móveis.

A pesquisa também apresenta abordagem \textbf{experimental}, uma vez que envolve a manipulação controlada de variáveis e a observação empírica dos resultados em ambientes reais. Conforme Lakatos e Marconi~\cite{lakatos2021}, o método experimental é adequado quando se pretende “testar hipóteses mediante o controle e a observação dos efeitos provocados por determinadas condições”.

\subsection{Objetivos da Abordagem}

A adoção de uma abordagem aplicada e experimental permite atender a objetivos técnicos e científicos específicos, entre os quais se destacam:

\begin{itemize}
  \item \textbf{Demonstrar a viabilidade técnica} de um sistema de monitoramento autônomo e totalmente embutido no ecossistema Flutter, sem dependência de ferramentas nativas externas (Android Developers~\cite{android2024});
  \item \textbf{Observar funcionalmente o impacto da coleta de métricas} sobre o desempenho da aplicação, verificando se o monitoramento permanece ativo sem comprometer a interação com o aplicativo de demonstração;
  \item \textbf{Verificar a coerência dos dados coletados}, confrontando-os qualitativamente com sinais observados no painel e com leituras complementares obtidas por ferramentas de referência, em especial comandos ADB documentados nos resultados~\cite{android2024};
  \item \textbf{Validar qualitativamente as recomendações heurísticas}, verificando se os alertas gerados são compatíveis com os gargalos induzidos nos cenários de teste;
  \item \textbf{Promover a reprodutibilidade científica}, disponibilizando publicamente o código-fonte e a documentação técnica do pacote, conforme recomendam Wohlin \textit{et al.}~\cite{wohlin2012}.
\end{itemize}

Esses objetivos orientam tanto o delineamento experimental quanto as etapas de desenvolvimento do sistema, assegurando coerência entre a fundamentação teórica, o projeto técnico e a avaliação empírica dos resultados.

\subsection{Justificativa da Escolha Metodológica}

A escolha de uma metodologia de natureza aplicada e abordagem experimental fundamenta-se na necessidade de validar empiricamente o comportamento de sistemas de monitoramento em contextos reais de execução.
Conforme destaca Yin~\cite{yin2015}, abordagens experimentais e estudos de caso são apropriados quando o pesquisador busca compreender relações de causa e efeito entre variáveis observáveis em ambientes complexos como é o caso de aplicações móveis em execução sobre diferentes plataformas e dispositivos.

No caso do \texttt{collector\_flutter}, a observação direta de variáveis como \textit{frame time}, reconstruções de widgets e uso de memória é indispensável para analisar se os sinais exibidos pelo painel acompanham os cenários executados.
Estudos empíricos recentes sobre desempenho móvel reforçam que tempo de execução, memória e energia devem ser avaliados de forma conjunta, pois cada indicador pode reagir de modo distinto às mudanças realizadas no software~\cite{rua2024largescale}. Por isso, a análise experimental é essencial para associar métricas observáveis a efeitos percebidos durante a execução do aplicativo.

Assim, a abordagem metodológica proposta combina:
\begin{itemize}
  \item fundamentos teóricos sobre desempenho, eficiência e sustentabilidade digital;
  \item execução controlada do pacote em ambiente Android;
  \item análise qualitativa dos resultados observados com apoio de ferramentas nativas de referência.
\end{itemize}

Essa integração entre teoria e prática sustenta a avaliação do sistema proposto e permite que ele seja reproduzido e aprimorado em diferentes contextos de desenvolvimento.



% ---------------------------------------------------------------------------------------------------------------------------------------------------------
\subsection{Ciclo de Vida da Pesquisa e do Desenvolvimento}
% ---------------------------------------------------------------------------------------------------------------------------------------------------------

O processo adotado neste trabalho segue um \textbf{ciclo de vida iterativo e incremental}, característico de projetos de engenharia de software com validação empírica, combinando desenvolvimento de pacote e experimentação controlada.
Esse ciclo foi estruturado com base em princípios de engenharia de software (Pressman e Maxim~\cite{pressman2020engenharia}; Sommerville~\cite{sommerville2019software}) e em diretrizes de experimentação em engenharia de software (Wohlin \textit{et al.}~\cite{wohlin2012}), garantindo rastreabilidade entre objetivos, decisões de projeto e resultados.

A Figura~\ref{fig:ciclo-vida} representa o ciclo metodológico adotado, organizado em fases com \textbf{entradas}, \textbf{atividades}, \textbf{artefatos gerados} e \textbf{critérios de término}.
O ciclo contempla dois momentos do trabalho: (i) \textbf{TCC 1}, voltado à fundamentação e ao protótipo inicial; e (ii) \textbf{TCC 2}, voltado à consolidação do pacote, aos testes funcionais e à análise qualitativa dos resultados.

\begin{figure}[H]
  \centering
  \includegraphics[width=0.95\textwidth]{ciclo_de_vida.png}
  \caption{Ciclo de vida metodológico adotado (iterativo e incremental).}
  \label{fig:ciclo-vida}
\end{figure}

A Figura~\ref{fig:metodologia-pesquisa} sintetiza a estratégia metodológica adotada no trabalho, articulando fundamentação teórica, levantamento de requisitos, construção do pacote e validação funcional. Essa organização evidencia que a avaliação realizada tem caráter qualitativo e funcional, sem análise estatística formal ou estudo externo de usabilidade.

\begin{figure}[H]
  \centering
  \includegraphics[width=0.95\textwidth]{metodologia_pesquisa.png}
  \caption{Estratégia metodológica da pesquisa, com foco em revisão bibliográfica, requisitos, implementação e validação funcional do pacote.}
  \label{fig:metodologia-pesquisa}
\end{figure}

\subsubsection{Visão Geral das Fases e Entregáveis}

O ciclo de vida foi operacionalizado em quatro fases principais, descritas a seguir. Para cada fase, definem-se claramente: (i) o objetivo, (ii) as atividades, (iii) os artefatos produzidos e (iv) o critério de conclusão.

\paragraph{Fase 1: Fundamentação e Revisão Bibliográfica}
\textbf{Objetivo:} construir o embasamento teórico sobre desempenho, profiling, métricas e eficiência energética em aplicações móveis, delimitando o problema e justificando a solução proposta.
\textbf{Atividades:} busca sistemática por palavras-chave, leitura e seleção de estudos, organização temática e consolidação de lacunas na literatura.
\textbf{Artefatos/saídas:} (i) revisão bibliográfica estruturada, (ii) definição do problema, (iii) objetivos e expectativas de validação.
\textbf{Critério de conclusão:} existência de um referencial teórico suficiente para sustentar as métricas monitoradas e as escolhas de projeto do pacote.

\paragraph{Fase 2: Requisitos, Escopo e Especificação}
\textbf{Objetivo:} transformar o problema de pesquisa em requisitos verificáveis e critérios de sucesso, descrevendo claramente o que o pacote deve fazer e quais limites deve respeitar.
\textbf{Atividades:} elicitação e análise de requisitos, definição de critérios de qualidade e operacionalização das métricas (o que medir e como medir).
\textbf{Artefatos/saídas:} (i) requisitos funcionais e não funcionais, (ii) critérios de avaliação, (iii) definição de metas (ex.: limite de \textit{overhead}).
\textbf{Critério de conclusão:} requisitos claros e mensuráveis, capazes de orientar implementação e avaliação experimental.

\paragraph{Fase 3: Projeto, Implementação e Integração do Pacote}
\textbf{Objetivo:} projetar e implementar o \texttt{collector\_flutter} de forma modular e testável, com integração a aplicações Flutter e capacidade de coleta/visualização local.
\textbf{Atividades:} definição da arquitetura (camadas e módulos), implementação incremental das funcionalidades, integração entre módulos e testes iniciais.
\textbf{Artefatos/saídas:} (i) pacote funcional, (ii) arquitetura documentada (diagramas e componentes), (iii) painel (\textit{dashboard}) integrado, (iv) registro e exportação de métricas.
\textbf{Critério de conclusão:} pacote executável, capaz de coletar e exibir métricas em tempo real em ambiente controlado.

\paragraph{Fase 4: Avaliação Experimental, Análise e Validação}
\textbf{Objetivo:} validar funcionalmente o pacote, observando se as métricas coletadas e as recomendações geradas são coerentes com os gargalos induzidos em cenários controlados.
\textbf{Atividades:} definição dos cenários experimentais, execução manual controlada, coleta automatizada de telemetria, inspeção dos resultados e confronto qualitativo com sinais obtidos por ADB.
\textbf{Artefatos/saídas:} (i) registros de telemetria, (ii) tabela de cenários, (iii) evidências visuais do painel, (iv) leituras complementares via ADB, (v) discussão de resultados e limitações.
\textbf{Critério de conclusão:} evidências suficientes para demonstrar a operacionalidade do pacote e a pertinência das heurísticas implementadas.

\subsubsection{Iterações e Rastreabilidade}

Embora as fases sejam apresentadas de forma sequencial, o ciclo adotado é \textbf{iterativo}: resultados parciais de testes e integração podem motivar ajustes em requisitos, arquitetura ou heurísticas, desde que tais mudanças sejam registradas e justificadas.
Para garantir rastreabilidade, cada funcionalidade do pacote foi associada a: (i) um requisito (RF/RNF), (ii) um módulo da arquitetura e (iii) um critério de avaliação experimental, permitindo ligação direta entre o que foi proposto, implementado e validado.

\subsubsection{Relação com o Cronograma do TCC 1 e TCC 2}

O ciclo metodológico foi distribuído ao longo do cronograma em dois macroperíodos:
\begin{itemize}
  \item \textbf{TCC 1:} foco nas Fases 1, 2 e início da Fase 3, resultando em fundamentação teórica consolidada e protótipo inicial funcional.
  \item \textbf{TCC 2:} foco na conclusão da Fase 3 e execução da Fase 4 em caráter funcional e qualitativo, resultando em análise dos cenários controlados e redação final da monografia.
\end{itemize}

Essa estrutura torna explícito o ciclo de vida adotado e define o que se espera em cada etapa do processo. A distribuição temporal completa das atividades de TCC 1 e TCC 2 está preservada no Apêndice~\ref{ap:cronogramas}, de modo a manter a rastreabilidade do planejamento sem deslocar o foco da metodologia.
